\chapter{Dynamical Systems -- Part 1}

\definecolor{boxbg}{HTML}{EEF3FB}
\definecolor{boxln}{HTML}{2E5EAA}
\definecolor{codebg}{HTML}{F5F5F5}

\lstset{basicstyle=\ttfamily\footnotesize,backgroundcolor=\color{codebg},
  frame=single,framesep=4pt,xleftmargin=4pt,breaklines=true,
  language=Python,showstringspaces=false,columns=fullflexible}

\begin{center}
{\Large\bfseries Dynamical Systems -- Part 1}\\[1pt]
{\small Exam cheat-sheet: \emph{continuous-time model, Euler / Runge-Kutta, equilibria \& stability}}
\end{center}
\hrule
\vspace{4pt}

A continuous-time component is the ``physics box''. To \textbf{specify} one, list:
\begin{itemize}
\item \textbf{Inputs} $I$: real-valued, type \texttt{real} or interval \texttt{real[L,U]}.
\item \textbf{Outputs} $O$ and \textbf{state} $S$: real-valued.
\item \textbf{Init}: the set of allowed start states $[\![Init]\!]$.
\item For each output $y$: an algebraic expression $h_y$ over $I\cup S$ \;($y(t)=h_y$).
\item For each state $x$: a derivative expression $f_x$ over $I\cup S$ \;($\tfrac{dx}{dt}=f_x$).
\end{itemize}
An \textbf{execution} given an input signal $I(t)$ is a state signal $S(t)$ + output signal $O(t)$ with:
\;(1) $S(0)\in[\![Init]\!]$;\;
(2) $y(t)=h_y(I(t),S(t))$;\;
(3) $\frac{d}{dt}x(t)=f_x(I(t),S(t))$.

\textit{Exam tip:} state variables get a \textbf{derivative} ($dx/dt=\dots$); outputs get a plain \textbf{equation} ($y=\dots$). Example car: state $v$ with $dv/dt=(F-kv)/m$, input force $F$.

\section{Does a solution exist / is it unique?}
Initial value problem (IVP): $\;\dfrac{dx}{dt}=f(x),\quad x(0)=x_0$.

\begin{keybox}
\textbf{Existence} (at least one solution): $f$ is \textbf{continuous}.\\
\textbf{Uniqueness} (exactly one solution) -- \emph{Picard--Lindel"of}: $f$ is \textbf{Lipschitz-continuous}.
\end{keybox}

\textbf{Lipschitz test:} $f$ is Lipschitz if there is a constant $K$ with
$\;|f(u)-f(v)|\le K\,|u-v|$ for all $u,v$ (a finite bound on how fast $f$ changes).

\subsection*{Quick classification (memorize these)}
\begin{itemize}
\item Linear, e.g. $(F-kv)/m$ \;$\Rightarrow$\; \textbf{Lipschitz} (always).
\item $x^2$ (any polynomial) \;$\Rightarrow$\; Lipschitz \textbf{only on a bounded domain}.
\item $x^{1/3}$ \;$\Rightarrow$\; \textbf{not} Lipschitz at $0$ $\Rightarrow$ IVP $dx/dt=x^{1/3},x(0)=0$ has \emph{multiple} solutions.
\item $f(x)=$ \texttt{if }$x=0$\texttt{ then }1\texttt{ else }0 \;$\Rightarrow$\; \textbf{not} continuous $\Rightarrow$ \emph{no} solution (jumps $\to$ hybrid systems).
\end{itemize}

\section{Euler method (the main computation)}
Numerically solve $\dfrac{dx}{dt}=f(x,t)$ from $x_0$ in steps of size $h$.

\begin{keybox}
\centering $\tilde x_{i+1}=\tilde x_i+h\cdot f(\tilde x_i,t_i)$,\qquad $t_{i+1}=t_i+h$,\qquad $\tilde x_0=x_0$.
\end{keybox}

\textbf{Recipe (fill the table):}
\begin{enumerate}
\item Write the starting row: $i=0,\;t_0,\;\tilde x_0$.
\item Compute slope $f(\tilde x_i,t_i)$ at the \emph{current} row.
\item New value $=$ old value $+\,h\times$ slope. Advance $t$ by $h$.
\item Repeat until $i=N$.
\end{enumerate}

\textbf{Worked example} ($\frac{dx}{dt}=t$, $x_0=0$, $t_0=0$, $h=1$, $N=5$; exact $x(t)=\tfrac12 t^2$):

\begin{center}\small
\begin{tabular}{cccc}
\toprule
$i$ & $t_i$ & $\tilde x_{i+1}=\tilde x_i+h\cdot f$ & exact $x(t_{i+1})$\\
\midrule
0 & 0 & $0+1\cdot0=0$ & 0.5\\
1 & 1 & $0+1\cdot1=1$ & 2\\
2 & 2 & $1+1\cdot2=3$ & 4.5\\
3 & 3 & $3+1\cdot3=6$ & 8\\
4 & 4 & $6+1\cdot4=10$ & 12.5\\
\bottomrule
\end{tabular}
\end{center}
\textit{Here $f$ depends only on $t$, so the slope is just $t_i$. If $f$ depends on $x$ (e.g. $dx/dt=-x$), plug in $\tilde x_i$ instead.}

\begin{lstlisting}
def euler(f, x0, t0, h, N):
    x, t = x0, t0
    for i in range(N + 1):
        print(i, round(t, 3), round(x, 4))
        x = x + h * f(x, t)   # Euler update
        t = t + h
euler(lambda x, t: t, 0, 0, 1, 5)   # dx/dt = t
# for dx/dt = -x  use:  lambda x, t: -x
\end{lstlisting}

\section{Runge--Kutta \& truncation error}
Euler uses only the slope at the \emph{start} of the interval; Runge--Kutta also looks ahead.

\textbf{2nd-order Runge--Kutta} (average of two slopes):
\[
k_1=f(\tilde x_i),\quad k_2=f(\tilde x_i+h\,k_1),\quad
\tilde x_{i+1}=\tilde x_i+h\,\tfrac{k_1+k_2}{2}.
\]
\textbf{4th-order Runge--Kutta} (most used): $\;\tilde x_{i+1}=\tilde x_i+\tfrac{h}{6}(k_1+2k_2+2k_3+k_4)$.

\begin{keybox}
\textbf{Global truncation error} (how far $\tilde x_i$ is from the true $x(t_i)$):\\
Euler: $\;|\tilde x_i-x(t_i)|\le C\,h$ \;(error $\propto h$).\qquad
RK4: $\;|\tilde x_i-x(t_i)|\le C\,h^4$ \;(error $\propto h^4$).\\
Smaller $h$ $\Rightarrow$ smaller error. Since $h^4<h$ for $0<h<1$, RK4 is more accurate than Euler.
\end{keybox}

\section{Equilibria}
\begin{keybox}
\centering A state $x_e$ is an \textbf{equilibrium} of $\dfrac{dx}{dt}=f(x)$ \;iff\; $f(x_e)=0$.
\end{keybox}
\textbf{Recipe:} set the derivative expression(s) to $0$ and solve for the state.
If you start exactly at $x_e$, the system stays there forever (constant signal).
\textit{Pendulum:} two equilibria -- hanging down ($\varphi=0$) and balanced straight up ($\varphi=-\pi$).

\section{Stability (Lyapunov)}
You perturb the system slightly away from $x_e$ -- what happens?
\begin{keybox}
\begin{tabular}{@{}ll@{}}
\textbf{Stable} & stays \emph{near} $x_e$ forever (small nudge $\to$ small wobble).\\
\textbf{Asymptotically stable} & stable \emph{and} returns to $x_e$ as $t\to\infty$ (stronger).\\
\textbf{Unstable} & drifts \emph{away} from $x_e$.
\end{tabular}
\end{keybox}
Formal: $x_e$ \emph{stable} if $\forall\varepsilon>0\,\exists\delta>0$: $|s-x_e|<\delta\Rightarrow |x[s](t)-x_e|<\varepsilon\ \forall t\ge0$.
\emph{Asymptotically stable} adds $\lim_{t\to\infty}x[s](t)=x_e$.\\
\textit{Pendulum intuition:} hanging down $=$ stable (frictionless) / asymptotically stable (with friction); balanced up $=$ unstable.

\vspace{3pt}
\hrule
\vspace{2pt}
{\footnotesize\textbf{One-line exam triage:} ``run Euler'' $\to$ table with $\tilde x_{i+1}=\tilde x_i+h f$; ``find equilibria'' $\to$ solve $f(x)=0$; ``unique solution?'' $\to$ check Lipschitz; ``error vs $h$'' $\to$ Euler $\propto h$, RK4 $\propto h^4$; ``classify $x_e$'' $\to$ stable / asymptotic / unstable.}
